% !TEX root = ../agr-paper.tex
% ---------------------------------------------------------------------
% From the thesis ablation study. ~1,400 words. THE PAPER'S HEADLINE
% SECTION since the September 2026 reframe.
%
% Delta is ablated-minus-reference, so POSITIVE MEANS REMOVING THE
% COMPONENT IMPROVED THE METRIC. Say that once, plainly, before the
% first number, or every figure below reads backwards.
%
% Report the planner asymmetry, never the pooled average -- pooling
% averages away the entire result. Only the WebQSP arm reaches
% significance; the CWQ arm is a DIRECTION, not an effect, and must not
% be called "marginally significant" (p = 0.088). The per-stratum
% breakdown (h1 improves on both datasets, h2 collapses on CWQ) is what
% turns the asymmetry from a curiosity into a mechanism; it was
% recomputed from the half-split records before being quoted.
%
% The three nulls get the same length as the positive result. That is
% the methodological argument, not an apology. The verification layer's
% null is the paper's own designed-around component producing nothing
% measurable, and the honest reading is the thesis's: it does not change
% how often AGR is right, it changes how often AGR is right for a reason
% it can exhibit -- with sec:limits-contribution bounding "exhibit".
%
% Power belongs HERE, not deferred to Discussion: n = 200 / 198, so a
% null is "no effect detected at this power". A reviewer will ask for
% the minimum detectable effect; it is computed in the text, and
% scripts/check_paper_numbers.py recomputes every figure in that
% subsection and pins each to its own sentence. Do not reword those
% sentences without re-running the check.
%
% What the design does NOT do is also said: one run per condition, one
% component at a time, no joint-removal condition. All three are the
% first things a reviewer asks for and they are in sec:outlook.
% ---------------------------------------------------------------------

\section{Component-Level Attribution}\label{sec:attribution}

Agentic systems are assemblies, and the literature reports them whole.
This section takes AGR apart. Four conditions each remove exactly one
component from the unmodified system, as \Cref{sec:removal} specifies,
and re-run the identical questions: no planner, no backtracking, no
claim verification, and embedding-only scoring ($\alpha = 1$, disabling
the model term in the edge scorer). Each condition is compared against
the reference on the same half-split, drawn as half of every hop stratum
rounded down, which gives $n = 200$ on WebQSP and $n = 198$ on
ComplexWebQuestions, with McNemar's test on the paired per-question
outcomes. The reference row is the unmodified system restricted to the
same questions, read from the main-run records rather than re-run, so it
is provably the system that produced \Cref{tab:main}. Each condition was
run once, and the components were removed one at a time; the
consequences of both choices are stated in \Cref{sec:power}.

\textbf{Throughout this section a delta is ablated minus reference, so a
positive number means removing the component \emph{improved} the
metric.} \Cref{tab:ablation} gives all four. The $p$-values are reported
as computed. A correction for multiple comparisons was not among the
pre-specified decisions, and applying one after seeing the results would
be a choice made on the results, so the four tests are read as
descriptive rather than as a family-wise claim.

\begin{table}[!tb]
  \begin{center}
        \caption{Component ablations. $\Delta$F1 is ablated minus reference:
      positive means the system got better without the component. The
      interval is a paired percentile bootstrap over questions
      ($10{,}000$ resamples); $p$ from McNemar's exact test on paired
      per-question correctness.}
    \label{tab:ablation}
    \scriptsize
    \setlength{\tabcolsep}{2.5pt}
    \begin{tabular}{|l|rlrr|rlrr|}
      \hline
      & \multicolumn{4}{c|}{\textbf{WebQSP} ($n=200$)}
      & \multicolumn{4}{c|}{\textbf{CWQ} ($n=198$)} \\
      \textbf{Removed} & $\Delta$F1 & $95\%$ CI & $p$ & Tokens
                       & $\Delta$F1 & $95\%$ CI & $p$ & Tokens \\
      \hline
      Planner        & $+0.083$ & $[+0.038, +0.130]$ & \textbf{0.006} & $-31\%$
                     & $-0.047$ & $[-0.108, +0.013]$ & 0.088 & $-21\%$ \\
      Backtracking   & $-0.013$ & $[-0.038, +0.010]$ & 0.727 & $-7\%$
                     & $0.000$  & $[-0.031, +0.030]$ & 1.000 & $-18\%$ \\
      Verification   & $+0.004$ & $[-0.006, +0.018]$ & 1.000 & $-2\%$
                     & $-0.009$ & $[-0.024, +0.000]$ & 1.000 & $-6\%$ \\
      Model scoring  & $-0.016$ & $[-0.060, +0.026]$ & 0.664 & $-25\%$
                     & $-0.008$ & $[-0.046, +0.030]$ & 0.481 & $-31\%$ \\
      \hline
    \end{tabular}
  \end{center}
\end{table}

\subsection{When Planning Hurts}\label{sec:planner}

One component produced a detectable effect, and it produced it in the
wrong direction. \textbf{Removing the planner improves WebQSP F1 by
$0.083$} ($p = 0.006$) while cutting token use by $31\%$. Twenty-one
questions were answered correctly only without the planner, against six
correct only with it. The disagreement is not marginal, and the cheaper
system is the more accurate one.

The mechanism is over-decomposition. WebQSP is predominantly a
single-hop benchmark: $256$ of its $400$ questions need one hop. Asked
to decompose a question that has no compositional structure, the planner
manufactures sub-objectives that are not really separable. The explorer
then scores candidate edges against a fragment of the question rather
than against the question, and it discards exactly the context that
would have identified the right relation. The planner prompt carries a
single-hop exemplar precisely to show that not decomposing is a valid
output, and that exemplar is evidently not enough. The per-stratum
breakdown locates the damage. On the WebQSP half-split, removing the
planner raises one-hop Hits@1 from $0.75$ to $0.85$ and two-hop Hits@1
from $0.84$ to $0.89$, so decomposition hurts simple and two-hop
questions alike on that benchmark.

On ComplexWebQuestions the direction reverses, as the mechanism
predicts. That benchmark is genuinely compositional, only $137$ of $400$
questions are single-hop, and removing the planner costs $0.047$ F1.
\textbf{That arm does not reach significance} ($p = 0.088$), and we
report it as a direction rather than an effect. We do not describe it as
marginally significant, because at this sample size that phrase would
assert more than the data carries. The strata show a coherent mechanism
rather than noise. Without the planner, one-hop Hits@1 \emph{improves}
from $0.49$ to $0.54$, two-hop collapses from $0.48$ to $0.34$, and
three-hop-plus moves from $0.62$ to $0.54$. That is exactly the pattern
in which decomposition helps genuine composition and hurts questions
that do not need it, and it is the pattern \Cref{tab:togstrata} shows
from the other side, where the agentic baseline beats AGR only on the
single-hop questions it finishes.

Pooling the two datasets would average a significant improvement against
a non-significant harm and report approximately nothing. The asymmetry
\emph{is} the result. Decomposition is not a general-purpose benefit. It
is a mechanism whose value depends on whether the question has
compositional structure to find. The actionable consequence follows
directly, and it is the recommendation we would most like to see
adopted: \textbf{gate decomposition on question structure} rather than
applying it unconditionally, predicting the number of hops from the
question and invoking the planner only above a threshold. On these
benchmarks that policy would capture the WebQSP improvement and the
ComplexWebQuestions capability at once.

\subsection{Three Components with No Detected Effect}\label{sec:nulls}

The remaining three conditions moved nothing detectable, and we report
them at the same length as the one that did. A field that reports only
what moved cannot attribute anything.

\textbf{Backtracking.} Removing it changes WebQSP F1 by $-0.013$ ($p =
0.727$) and ComplexWebQuestions F1 by $0.000$ ($p = 1.000$), while
saving $7\%$ and $18\%$ of tokens. The mechanism fires rarely. The
reference condition records $38$ and $108$ backtracks over the $398$
paired questions, and on those same runs the meter refuses a further
attempt on $28$ of them, $7.0\%$. A component this seldom exercised
cannot show a large effect. One implementation detail bounds the null
further, and we record it rather than smooth it over. The backtracker
bans the edges of the most recent explorer pass but restores the
highest-scoring snapshot, which is often older, so everything expanded
in between stays unbanned (\Cref{sec:statemachine}). Over the main runs,
$53$ explorer passes re-expanded an anchor-and-relation set the same
question had already expanded, and at least $65$ of the $262$ backtracks
restored a depth other than the most recent snapshot's. The null
therefore describes backtracking as implemented, not backtracking as
specified. The accurate summary is that the search rarely needs to undo
a decision, not that undoing is worthless.

\textbf{Model scoring.} Setting $\alpha = 1$ removes the language model
from edge selection entirely, leaving pure embedding similarity. F1
falls by $0.016$ and $0.008$, neither significant, while tokens drop
$25\%$ and $31\%$. The model term is the single most expensive part of
the navigation loop, and this experiment cannot establish that it earns
its cost. The development-set sweep that froze $\alpha$ points the same
way at a smaller scale: at $\tau = 0.20$ it scored $61$, $62$, $64$, and
$59$ hits of $80$ at $\alpha = 0.3$, $0.5$, $0.7$, and $1.0$, so
dropping the model term cost five hits there, before the entity-filter
correction that preceded the test runs. One confound is recorded rather
than smoothed over. At $\alpha = 1$ the scorer takes an early return
that leaves the backtracking trigger's signal maximum computed over a
smaller candidate set, so this condition's elevated backtrack count is
not attributable to the scoring change alone.

\textbf{Claim verification.} Removing the checking routes of the layer
AGR was designed around, while leaving the drafting call unchanged,
changes WebQSP F1 by $+0.004$ and ComplexWebQuestions F1 by $-0.009$, $p
= 1.000$ on both. Across the half-splits --- $200$ questions on WebQSP
and $198$ on ComplexWebQuestions --- the two conditions disagreed on
exactly one question each. \textbf{Claim verification does not
measurably change how often AGR is right.}

The result is expected rather than disappointing, for the reason
\Cref{sec:verification} gave before any of these numbers appeared. On
most questions the layer is a pass-through certifying a draft that was
already grounded. On test it asks for a repair on $52$ of $800$
questions, exploration actually resumes on $28$, and $13$ end grounded
that did not start so. A mechanism that alters the output on a few per
cent of questions cannot move an aggregate accuracy metric. An ablation
is the right instrument for establishing that, rather than assuming it
either way.

What the layer delivers is therefore not an accuracy gain but the output
contract: an answer paired with the traversed triples that support it,
and an entity list that is a deterministic function of the verdicts
rather than a second generation. That is an auditability claim, not an
accuracy claim, and \Cref{sec:limits-contribution} states how far it
reaches: which route accepts most claims is not recorded, and the
supporting triples are counted rather than kept. \Cref{sec:groundedness}
has already shown that the layer cannot claim credit for groundedness.

\subsection{Statistical Power}\label{sec:power}

These are nulls at low power, not demonstrations of absence, and the
arithmetic is worth stating exactly. McNemar's exact test sees only the
\emph{discordant} pairs, the questions on which the two conditions
disagree, and the null conditions produced very few of them.

Removing claim verification changed the outcome on exactly one question
per dataset. With a single discordant pair no split whatsoever can reach
significance. So that row is not a test that failed to reject. It is the
two conditions agreeing on $396$ of the $398$ paired questions, which is
a statement about how rarely the layer alters an answer rather than
about statistical power. Backtracking produced $8$ and $11$ discordant
pairs, and model scoring $21$ and $18$. At those counts the smallest
imbalance the exact test could have called significant is $8$ and $9$
questions for backtracking, meaning all eight would have had to fall the
same way, and $11$ and $10$ for model scoring, which is between $4.0$
and $5.5$ points of accuracy. Differences below that were undetectable
at any power, whatever their true size.

For calibration: detecting a true $2{:}1$ discordance ratio at $80\%$
power requires roughly $72$ discordant pairs, and the largest of these
conditions produced $21$. At that count the smallest ratio detectable at
$80\%$ power is nearer $4{:}1$. A true $2{:}1$ effect would have been
caught about a quarter of the time. \textbf{Three of four components
show no effect at this power}; none is shown to have no effect.

One more count bounds what any single removal could have done. The union
of the four discordant sets holds $42$ of the $200$ WebQSP questions and
$55$ of the $198$ on ComplexWebQuestions. On the remaining $158$ and
$143$, no condition changed the outcome. The intervals in
\Cref{tab:ablation} say the same thing in F1: the planner's WebQSP
interval is the only one of the eight that does not contain zero.

Two features of the design bound these numbers further. Each condition
is one run against a hosted model whose temperature-zero decoding is
greedy rather than deterministic, so a delta of a few hundredths sits
inside the run-to-run variance that \Cref{sec:limits-comparison}
reports, and none of the three nulls was replicated. And the components
were removed one at a time. Three nulls in one-at-a-time removal are
consistent with the three components being individually dispensable and
jointly necessary, and equally with a system that has none of them
matching this one at a fraction of the tokens. The condition that
separates those readings, all three removed together, was not run.
\Cref{sec:outlook} names it, with the replication, as the next
experiment.
